One may therefore suppose $Z$ also affine, $Z=\Spec(B)$, where $B$ is a free $A$-module. Let $J$ be a subset of $B$ defining the subscheme $Y$ of $Z$, and let $K$ be the ideal in $A$ generated by the $u_i(J)\subseteq A$, where the $u_i:B\to A$ are the coordinate forms with respect to the chosen basis. One verifies immediately that $T=\mathcal V(K)=\Spec(A/K)$ satisfies the desired condition, which completes the proof.
\begin{examples}{6.5}\label{VIII.6.5}
Let us give important examples of functors that reduce to functors $\prod_{Z/S}Y/Z$ of the type considered in 6.4 and for which it will be useful in what follows to have criteria of representability. Denote by $S$ a prescheme, by $X,Y,Z$, etc., preschemes over $S$.

a) Let us give ourselves an $S$-morphism
\begin{equation}\tag{x}
q:X\to\uHom_S(Y,Z),
\end{equation}
(``$X$ acts on $Y$, with values in $Z$''), i.e. a morphism
\begin{equation}\tag{xx}
r:X\times_S Y\to Z.
\end{equation}
Consider a subprescheme $Z'$ of $Z$, whence a monomorphism
\[
\uHom_S(Y,Z')\to\uHom_S(Y,Z)
\]
which makes the first functor a subfunctor of the second; let $X'$ be the inverse image of this subfunctor by (x). It is the subfunctor of $X$ such that $X'(T)$ is the set of $x\in X(T)$ such that $q(x):Y_T\to Z_T$ factors through $Z'_T$. This functor $X'$ may be described as follows: put $P=X\times_S Y$, let $P'$ be the inverse image of $Z'$ by $r:P\to Z$; then one has an evident isomorphism
\begin{equation}\tag{xxx}
X'\simeq\prod_{P/X}P'/P.
\end{equation}
One therefore obtains: \emph{if $Y$ is essentially free over $S$ and $Z'$ closed in $Z$, the subfunctor $X'$ of $X$ is representable by a closed subprescheme of $X$.}

b) Let us give ourselves two ways of making $X$ act on $Y$ with values in $Z$, i.e. two morphisms
\[
\xymatrix{q_1,q_2:X\ar@<.5ex>[r]\ar@<-.5ex>[r]&\uHom_S(Y,Z),}
\]
and put $X'=\Ker(q_1,q_2)$: it is the subfunctor of $X$ such that $X'(T)$ is the set of $x\in X(T)$ such that the two morphisms $q_1(x),q_2(x):Y_T\rightrightarrows Z_T$ are equal. Now giving $q_1,q_2$ is equivalent to giving a morphism
\[
q:X\to\uHom_S(Y,Z\times_S Z),
\]
or again a morphism $r:X\times_S Y\to Z\times_S Z$; put then $U=Z\times_S Z$, let $U'$ be the diagonal subprescheme of $Z\times_S Z$; then $X'$ is precisely the inverse image of the subfunctor $\uHom_S(Y,U')\to\uHom_S(Y,U)$ by $q$, and so can be put in the form (xxx), with $P=X\times_S Y$, and $P'=$ inverse image of the diagonal by $r$, i.e. kernel of $X\times_S Y\mathrel{\substack{\xrightarrow{r_1}\\[-.8ex]\xrightarrow{r_2}}}Z$. One is therefore under the conditions of (a).

One thus sees that: \emph{if $Y$ is essentially free over $S$ and $Z$ separated over $S$, then the subfunctor $X'$ of $X$ is representable by a closed subprescheme of $X$.}

c) Let us give ourselves a morphism
\[
q:X\to\uHom_S(Y,Y),
\]
i.e. ``$X$ acts on $Y$''. Let $X'$ be the ``kernel'' of this morphism, i.e. the subfunctor $X'$ of $X$ such that $X'(T)$ is the set of $x\in X(T)$ such that $q(x):Y_T\to Y_T$ is the identity. This functor is covered by b), as one sees by introducing a second homomorphism
\[
q':X\to\uHom_S(Y,Y)
\]
``by making $X$ act trivially on $Y$''. Thus: \emph{if $Y$ is essentially free over $S$ and separated over $S$, the kernel subfunctor of $q$ is representable by a closed subprescheme of $X$.}

d) Under the conditions of c), consider the subfunctor $Y'$ of $Y$ ``of invariants under $X$'', so $Y'(T)$ is the set of $y\in Y(T)$ such that the corresponding morphism $\overline q(y):X_T\to Y_T$ is ``the constant $T$-morphism with value $y$''. Introducing $q'$ as in c), and the homomorphisms corresponding to $q$ and $q'$:
\[
\xymatrix{\overline q,\overline{q'}:Y\ar@<.5ex>[r]\ar@<-.5ex>[r]&\uHom_S(X,Y),}
\]
one sees that $Y'$ is precisely $\Ker(\overline q,\overline{q'})$, and is therefore again covered by b) (with the roles of $X,Y$ reversed and $Z=Y$).

Consequently, \emph{if $X$ is essentially free over $S$, $Y$ separated over $S$, then the subfunctor $Y'$ of $Y$ of invariants under $X$ is representable by a closed subprescheme of $Y$.}

e) Constructions of the type made explicit in the preceding examples are especially frequent in group theory. Thus, when $G$ is an $S$-group prescheme acting on the $S$-prescheme $X$:
\[
q:G\to\underline{\Aut}_S(X),
\]
the kernel of $q$ (``the subgroup of $G$ acting trivially'') is a closed subscheme of $G$ provided that $X$ is essentially free and separated over $S$ (example c)), and the subobject $X^G$ of invariants is a closed subprescheme of $X$, provided that $G$ is essentially free over $S$ and $X$ separated over $S$\nde{$S_X$ has been corrected to $S$.} (example d)).

Let $Y,Z$ be subpreschemes of $X$; consider the subfunctor $\underline{\operatorname{Transp}}_G(Y,Z)$ of $G$ (``transporter of $Y$ into $Z$'') whose points with values in a $T$ over $S$ are the $g\in G(T)$ such that the corresponding automorphism of $X_T$ satisfies $g(Y_T)\subseteq Z_T$, i.e. induces a morphism $Y_T\to X_T$ factoring as $Y_T\to Z_T$. Thus: \emph{if $Y$ is essentially free over $S$ and $Z$ closed in $X$, then $\underline{\operatorname{Transp}}_G(Y,Z)$ is a closed subprescheme of $G$} (example a)).

One may also consider the \emph{strict transporter of $Y$ into $Z$},\nde{Denoted by $\underline{\operatorname{Transp}}^{\mathrm{str}}_G(Y,Z)$.} whose points with values in a $T$ over $S$ are the $g\in G(T)$ such that $g(Y_T)=Z_T$, which is precisely $\underline{\operatorname{Transp}}_G(Y,Z)\cap\sigma(\underline{\operatorname{Transp}}_G(Z,Y))$, where $\sigma$ is the symmetry of $G$. Consequently, \emph{if $Y$ and $Z$ are essentially free over $S$ and closed in $X$, the strict transporter of $Y$ into $Z$ is a closed subprescheme of $G$.}

An important case is that where $X=G$, $G$ acting on itself by inner automorphisms. If $H$ is a subprescheme of $G$, the strict transporter of $H$ into $H$ is also called the \emph{normalizer of $H$ in $G$}, and denoted by $\uNorm_G H$. Thus: \emph{if $H$ is a closed subgroup prescheme of $G$, essentially free over $S$, then $\uNorm_G H$ is representable by a closed subgroup prescheme of $G$.}

Finally, let $Z$ be a subprescheme of $G$; then its \emph{centralizer} $\uCentr_G(Z)$ in $G$ is the group subfunctor of $G$ defined by the procedure of d), when one considers that ``$Z$ acts on $G$'' by the operations induced by those of $G$; thus \emph{if $Z$ is essentially free over $S$ and $G$ separated over $S$, $\uCentr_G(Z)$ is a closed subgroup prescheme of $G$.} In particular, \emph{if $G$ is essentially free and separated over $S$, then the center $C$ of $G$, which is precisely $\uCentr_G(G)$, is a closed subgroup prescheme of $G$.}
\end{examples}
When $S$ is the spectrum of a field, 6.3 b) shows that in examples a) to e) above the ``essentially free'' conditions are automatically satisfied; only separation conditions remain. Recalling that a group prescheme over a field is necessarily separated, one finds, for example:
% typo?: reference 6.3 b) retained as printed.
\begin{corollary}{6.7}\label{VIII.6.7}
\nde{The original numbering has been preserved: there is no 6.6.} Let $G$ be a group prescheme over a field $k$. Then:
\begin{itemize}
\item For every subprescheme $Z$ of $G$, the centralizer of $Z$ in $G$ is a closed subgroup prescheme of $G$; this is in particular the case for the center $\uCentr_G(G)$ of $G$.
\item More generally, if $u,v:X\to G$ are morphisms of preschemes, $\underline{\operatorname{Transp}}_G(u,v)$ is representable by a closed subprescheme of $G$.
\item For two subpreschemes $Y,Z$ of $G$, with $Z$ closed, $\underline{\operatorname{Transp}}_G(Y,Z)$ is a closed subprescheme of $G$. If $Y$ is also closed, one has the same conclusion for $\underline{\operatorname{Transp}}^{\mathrm{str}}_G(Y,Z)$.
\item For every subgroup prescheme\nde{Indeed, over a field $k$, every subgroup prescheme of $G$ is closed, cf. VI$_{\mathrm A}$, 0.6.1.} $H$ of $G$, the normalizer $\uNorm_G(H)$ is a closed subgroup scheme of $G$.
\end{itemize}
\end{corollary}
\begin{remark}{6.8}\label{VIII.6.8}
\nde{The original has been expanded in what follows; in particular lemma 6.8.1 has been added.} Let $A$ be a commutative ring, $M$ an $A$-module, $M^\vee=\Hom_A(M,A)$; endow the $A$-module $\End_A(M)$ with the topology of discrete pointwise convergence, i.e. a neighborhood basis of 0 consists of the following $A$-submodules, where $n\in\NN$ and $m_1,\dots,m_n\in M$:
\[
K(m_1,\dots,m_n)=\{u\in\End_A(M)\mid u(m_i)=0\quad\text{for }i=1,\dots,n\}.
\]
One says that $M$ is a \emph{quasi-free $A$-module} if the image of the canonical morphism $\Theta:M\otimes_A M^\vee\to\End_A(M)$ contains $\mathrm{id}_M$ in its closure, i.e. if the following condition is satisfied:
\begin{equation}\tag{*}
\begin{gathered}
\text{for every }m_1,\dots,m_n\in M,\text{ there exist }x_1,\dots,x_r\in M\text{ and }f_1,\dots,f_r\in M^\vee\\
\text{such that }m_i=\Theta\left(\sum_{s=1}^r x_s\otimes f_s\right)(m_i)=\sum_{s=1}^r f_s(m_i)x_s\text{ for }i=1,\dots,n.
\end{gathered}
\end{equation}
(In this case, $\Im\Theta$ is dense in $\End_A(M)$ since for every $u\in\End_A(M)$ one has $u(m_i)=\sum_{s=1}^r f_s(m_i)u(x_s)=\Theta(\sum_{s=1}^r u(x_s)\otimes f_s)(m_i)$.)

Let us first note that this property is stable under base change. Indeed, let $\phi:A\to A'$ be a ring morphism, $M'=M\otimes_A A'$, and $m'_1,\dots,m'_n\in M'$; then $m'_i=\sum_j m_{ij}\otimes b_{ij}$ ($m_{ij}\in M$, $b_{ij}\in A'$); by hypothesis, there exist $x_1,\dots,x_r\in M$ and $f_1,\dots,f_r\in\Hom_A(M,A)$ such that $m_{ij}=\sum_s x_sf_s(m_{ij})$ for every $i,j$. Denote by $\phi\circ f_s$ the image of $f_s$ in $\Hom_A(M,A')=\Hom_{A'}(M,A')$; then for every $i=1,\dots,n$ one has:
\[
\left(\sum_s x_s\otimes\phi\circ f_s\right)(m'_i)=\sum_{s,j}x_s\otimes\phi(f_s(m_{ij}))b_{ij}=\sum_j\left(\sum_s x_sf_s(m_{ij})\right)\otimes b_{ij}=m'_i,
\]
which proves that $M'$ is quasi-free over $A'$.
% typo?: Hom_{A'}(M,A') retained as printed.
Let us also note that every projective $A$-module $P$ is quasi-free (there exist $A$-morphisms $A^{(I)}\xrightarrow{\pi}P\xrightarrow{\tau}A^{(I)}$ such that $\pi\circ\tau=\mathrm{id}_P$; denote by $(e_i)$ the canonical basis of $A^{(I)}$ and by $f_i$ the linear form $e_i^*\circ\tau$ on $P$; if $m_1,\dots,m_n\in P$, there exists a finite subset $J$ of $I$ such that $m_k=\sum_{i\in J}f_i(m_k)\pi(e_i)$ for $k=1,\dots,n$).

Then theorem 6.4 remains valid on replacing, in the statement of definition 6.1, the word ``free'' by ``projective'' or, more generally, by ``quasi-free''. Indeed, proceeding as in the proof of 6.4, one is reduced to proving the
\end{remark}
\begin{lemma}{6.8.1}\label{VIII.6.8.1}
Let $M$ be a quasi-free $A$-module, $N$ a submodule, $F$ the covariant functor $(A\text{-algebras})\to\Ens$ such that $F(B)=\{\mathrm{pt}\}$ if $M_B=(M/N)_B$, and $F(B)=\varnothing$ otherwise. Then there exists an ideal $K$ of $A$ such that $F(B)=\{\mathrm{pt}\}$ if and only if the morphism $A\to B$ factors through $A/K$, i.e. one has an isomorphism functorial in $B$:
\[
F(B)\simeq\Hom_{A\text{-}\mathrm{alg.}}(A/K,B).
\]
\end{lemma}
\begin{proof}
Let $(n_j)$ be a system of generators of $N$, and let $F_j$ be the subfunctor of the final functor $\underline e$ ($\underline e(B)=\{\mathrm{pt}\}$ for every $B$) corresponding to the submodule $An_j$. One has $F(B)=\{\mathrm{pt}\}$ if and only if the image of each $n_j$ in $M_B$ is zero, so $F$ is the intersection of the functors $F_j$. This reduces us to the case where $N$ is generated by an element $n$.

Let $\phi:A\to B$ be an $A$-algebra; if the image $n\otimes1$ of $n$ in $M_B$ is zero, then for every $f\in M^\vee=\Hom_A(M,A)$ one has $0=f(n)\otimes1=\phi(f(n))$. On the other hand, since $M$ is quasi-free, there exist $x_1,\dots,x_r\in M$ and $f_1,\dots,f_r\in M^\vee$ such that $n=\sum_s f_s(n)x_s$, whence $n\otimes1=\sum_s x_s\otimes\phi(f_s(n))$. It follows that $n\otimes1=0$ if and only if $\phi$ factors through $A/K$, where $K$ is the ideal generated by the $f_s(n)$ for $s=1,\dots,r$. This proves the lemma.
\end{proof}

\sgasection{7}{Appendix: On monomorphisms of group preschemes}
The result proved in the present \No\ is unnecessary for the rest of the seminar, except X 8.8 and XV and XVI. It relies essentially on the existence theorem for quotient groups of Expos\'e VI$_{\mathrm A}$.\nde{Cf. Exp. VI$_{\mathrm A}$, Theorems 3.2 and 3.3.2.}

Corollary 5.7 leads one to ask under what conditions one can assert that a monomorphism $u:G\to H$ of $S$-groups is an immersion, or even a closed immersion.

We have seen in VI$_{\mathrm B}$ 1.4.2 that this is so if $S$ is the spectrum of a field, provided $G$ is of finite type over $k$ and $H$ locally of finite type over $k$. One easily concludes that the same result remains valid if one merely supposes $S$ artinian.\nde{Take account of the additions made in VI$_{\mathrm B}$\dots}

On the other hand, it is easy to give examples of bijective monomorphisms which are not immersions, $S$ being, for example, the affine line over a field or the spectrum of a discrete valuation ring. One will take, for example, $H=(\ZZ/2\ZZ)_S$, $G=G_1\times_S G_2$, where $G_1$ is the open subgroup of $H$ complementary to the closed point $x$ distinct from 0 of the fiber $H_s\simeq\ZZ/2\ZZ$ (where $s$ denotes a fixed closed point of $S$), and $G_2$ the closed subscheme of $H$ which is the sum of the subscheme reduced to the unit section and of the reduced closed subscheme defined by the closed point $x$. One easily verifies that $G_2$ is indeed stable under multiplication in $H$, and hence is a group scheme. The immersions $G_i\to H$ ($i=1,2$) then define a homomorphism of $S$-groups $G=G_1\times G_2\to H$, which is evidently a bijective monomorphism (and in addition a local immersion), but is not an immersion. (One verifies that $G$ and $H$ are reduced, $G$ having three disjoint irreducible components, whereas $H$ has only two irreducible components.) One will note that $G_1\to H$ also gives an example of an open subgroup $G$ of $H$ which is not closed (contrary to what occurs for algebraic groups over a field). The theory of degeneration of elliptic curves supplies other examples of this latter phenomenon, with $H$ in addition smooth over $S$ of relative dimension 1, and $G$ with connected fibers.
% typo: "deux composante irréductibles" -> "two irreducible components".
It is possible,\footnote{In fact, M. Raynaud has constructed a counterexample, with $G$ smooth with connected fibers, cf. XVI 1.1 c). If one does not suppose $G$ to have connected fibers, one may take $S=\Spec(\ZZ_2)$, $G=(\ZZ/2\ZZ)_S$ with the non-neutral closed point removed, $H=(\bmu_2)_S$.} on the other hand, that as soon as one supposes $G$ flat over $S$, and (say) $G$ and $H$ of finite presentation over $S$, a monomorphism $u:G\to H$ of $S$-groups is automatically an immersion. We shall prove a result of this kind, under additional hypotheses.

Let us first note that one may suppose $S$ affine, and (thanks to the finite presentation hypothesis on $G$ and $H$, which allows one to reduce to the case of the spectrum of a ring of finite type over $\ZZ$\nde{Cf. EGA IV$_3$, \S8, and Exp. VI$_{\mathrm B}$, \S10.}) $S$ noetherian. Then $G$ and $H$ are noetherian. Saying that $u:G\to H$ is a closed immersion (resp. an immersion) then amounts to saying that $u$ is a monomorphism (which is true by hypothesis) and that $u$ is proper (resp. and that $u$ is proper at every point of $u(G)$).\nde{Cf. EGA IV$_3$, 8.11.5 and 15.7.1.} The valuative criterion of properness assures us that it suffices to check that for every base change $S'\to S$, with $S'$ the spectrum of a discrete valuation ring, complete if one wishes, the morphism $u':G'\to H'$ has the same properness property. (The case of local properness was forgotten in EGA II 7.3, and will appear as an afterthought in EGA IV.\footnote{Cf. EGA IV$_3$, 15.7.}) This therefore reduces us to the case where \emph{$S$ is itself the spectrum of a complete discrete valuation ring}, --- provided that the additional hypotheses on $G,H,u$ which we shall be led to formulate are stable under base change.

Let then $s$ (resp. $s'$) be the closed point (resp. the generic point) of $S$. Then the homomorphisms induced on the fibers
\[
u_s:G_s\to H_s\quad\text{and}\quad u_{s'}:G_{s'}\to H_{s'}
\]
are closed immersions, since they are monomorphisms of group schemes of finite type over fields (VI$_{\mathrm B}$ 1.4.2). We may therefore identify $G_{s'}$ with a closed subscheme of $H_{s'}$. Now one has the following result:
\begin{lemma}{7.1}\label{VIII.7.1}
Let $S$ be the spectrum of a discrete valuation ring, $s'$ its generic point, $H$ an $S$-prescheme, $L'$ a closed subprescheme of the generic fiber $H_{s'}$, so that $L'$ is also a subprescheme of $H$.

Then the schematic closure $\overline{L'}$ in $H$ (i.e. the smallest closed subprescheme of $H$ containing $L'$, cf. EGA I 9.5) exists and is also the unique closed subprescheme of $H$, flat over $S$, whose generic fiber is $L'$. Moreover, formation of $L$ is functorial with respect to a variable pair $(H,L')$ and commutes with formation of cartesian products over $S$.

In particular, if $H$ is an $S$-group and $L'$ a subgroup of $H_{s'}$, then $L$ is a subgroup of $H$.
\end{lemma}
The proof is immediate and left to the reader.\footnote{Cf. EGA IV$_2$, 2.8.} Applying this to the situation $H,G_{s'}$, one sees that the monomorphism $u:G\to H$ factors as $G\to L\to H$, where $L$ is a subgroup of $H$ which is a closed subprescheme, flat over $S$, and where $G\to L$ induces an isomorphism on the generic fibers. Then $u:G\to H$ is an immersion (resp. a closed immersion) if and only if $u':G\to L$ is so. This therefore reduces us to the case where $H$ is flat over $S$ and $u_{s'}:G_{s'}\to H_{s'}$ an isomorphism (provided that the additional hypotheses which we shall have to formulate on $G,H,u$ are respected when one replaces $H$ by a closed subgroup prescheme). As $H$ is then the schematic closure of $H_{s'}$, if $u$ is an immersion, $u(G)$ will be a subprescheme\nde{The word ``closed'' has been deleted.} of $H$ containing $H_{s'}$, so its schematic closure will also be $H$, and consequently $u(G)$ will be an open subprescheme of $H$. Consequently, we shall in fact have to prove that $u$ is an open immersion (resp. an isomorphism if we wished to establish that $u$ is a closed immersion). As $G$ and $H$ are flat over $S$, it amounts to the same thing to say that the morphisms induced on the fibers are open immersions, resp. isomorphisms (cf. SGA 1, I 5.7), and since this is already so for $u_{s'}$, one is reduced to proving that $u_s$ is an open immersion, resp. an isomorphism.
